% !PRINTED none
\begin{center}
\gesperrt{Sechzehnter Abschnitt.}

\vspace{1em}

\textbf{Klassenzahl in quadratischen K\"orpern.}

\vspace{0.5em}
\rule{4cm}{0.4pt}
\end{center}

\begin{center}
§ 110. \gesperrt{Fundamentale Einheiten in den Ordnungen.}
\end{center}

Die Theorie der Einheiten, die wir allgemein in Bd. II, § 191
auseinandergesetzt haben, nimmt für die quadratischen Körper
eine einfachere Gestalt an, mu\ss{} aber andererseits erweitert und
für die \gesperrt{Ordnungen} ausgebildet werden.

Die ganzen Zahlen der Ordnung $O$ mit der Diskriminante $D$
sind nach § 96, (6) in der Form $x_1 + x_2\theta$ enthalten, worin
\[
\begin{array}{ll}
\theta = \tfrac{1}{2}\sqrt{D}, & \text{wenn } D \equiv 0 \ (\mathrm{mod}\ 4),\\[4pt]
\theta = \dfrac{1+\sqrt{D}}{2}, & \text{„}\quad D \equiv 1 \ (\mathrm{mod}\ 4),
\end{array}
\]
und $x_1, x_2$ ganze rationale Zahlen sind.

Eine solche Zahl $\varepsilon$ ist eine Einheit, wenn

\weq{1}{N(\varepsilon) = (x_1+x_2\theta)(x_1+x_2\theta') = \pm 1}

ist. Hierfür ergibt sich die Bedingung

\weq{2}{x_1^2 - x_2^2 D = \pm 4.}

Ist $D$ negativ, so kann hier nur das positive Zeichen stehen,
und es gibt im allgemeinen nur die zwei Lösungen $x = \pm 2$,
$y = 0$ % [sic] Gedruckt sind x und y statt x_1 und x_2. und in den beiden Ausnahmefällen die vier oder sechs
Lösungen:
\[
\begin{array}{lll}
D = -4, & x_1 = \pm 2, & x_2 = 0,\\
 & x_1 = 0, & x_2 = \pm 1.\\
D = -3, & x_1 = \pm 2, & x_2 = 0,\\
 & x_1 = \pm 1, & x_2 = \pm 1.
\end{array}
\]
